\section{Технический проект}

\subsection{Общая характеристика организации решения задачи}

Необходимо спроектировать и разработать оптимизирующий компилятор для языка Common Lisp, который способен транслировать исходное S-выражение в байт-код для целевой виртуальной машины. Компилятор должен поддерживать основные конструкции, а также применять ряд оптимизаций для повышения эффективности генерируемого кода.

Компилятор представляет собой программную систему, которая принимает на вход S-выражение и флаги оптимизации, преобразует выражение в промежуточное представление, применяет оптимизации и генерирует байт-код.

\subsection{Архитектура оптимизирующего компилятора Common Lisp}

\subsubsection{Общая архитектура компилятора}

Компилятор должен быть представлен в виде набора модулей, каждый из которых отвечает за ту или иную фазу компиляции, кроме основного модуля compiler. Должны быть определены следующие модули:

\begin{enumerate}
\item compiler. Включает в себя функцию (compile expr optimization-flags), которая в свою очередь поочерёдно вызывает функции из других модулей, передавая результат очередной функции на вход следующей.
\item analyzer. Состоит из функции (analyze expr), которая принимает на вход исходное S-выражение expr и производит его статический анализ, а также внутренних вспомогательных функций для анализа отдельных подвыражений. Функция analyze возвращает упрощённую промежуточную форму искомого S-выражения в случае успеха и сообщение об ошибке в противном случае.
\item optimizer. Состоит из функции (optimize tree flags), которая принимает на вход промежуточную форму S-выражения tree, полученную на выходе из функции analyze, и преобразует данную форму, применяя к ней оптимизации, указанные в качестве списка флагов оптимизаций flags, а также вспомогательных функций, отвечающих за конкретные оптимизации. Функция opimize возвращает промежуточную форму с применёнными к ней оптимизациями, которая может остаться без изменений, если ни одна из указанных оптимизаций неприменима к данной промежуточной форме, либо если в качестве списка флагов оптимизаций на вход функции передан пустой список.
\item generator. Состоит из функции (generate tree), которая принимает на вход оптимизированную промежуточную форму S-выражения tree, полученную на выходе из функции opimize, и на её основе генерирует код ассемблера, состоящий из мнемоник инструкций, их операторов и меток для инструкций перехода, а также вспомогательных функций для генерации отдельных элементов промежуточной формы. Функция generate возвращает массив мнемоник инструкций, их операндов и меток.
\item assembler. Состоит из функции (assemble code), которая принимает на вход ассемблер код -- массив мнемоник инструкций, операндов инструкций и меток для перехода, полученный на выходе из функции generate, и переводит его в байт-код для виртуальной машины, а также других вспомогательных функций. Функция assemble возвращает массив байтов, составляющих искомый байт-код, список констант, используемых в программе и размер памяти для глобальных переменных в случае успешного ассемблирования, а в случае ошибки -- прекращает своё выполнение с сообщением об ошибке.
\item vm. Представляет собой целевую виртуальную машину и состоит из основной функции (vm-run bytecode), которая выполняет байт-код, полученный модулем assembler, вычисляя исходную программу, а также дополнительных вспомогательных функций. Функция vm-run возвращает на выходе итоговое вычисленное значение исходного скомпилированного S-выражения, либо заканчивается ошибкой выполнения и возвращает в таком случае сообщение об ошибке.
\end{enumerate}

\subsubsection{Фаза 1. Статический анализ}

На данном этапе компиляции входное S-выражение преобразуется в другое абстрактное синтаксическое дерево, состоящее из простых операций и сохраняющее смысл исходного выражения.

\paragraph{Лексическое окружение}

Для статического разрешения переменных используется лексическое окружение. Оно представляет собой набор кадров, каждый из которых содержит список связанных переменных. В начале компиляции окружение пустое. При анализе определения функции окружение расширяется новым кадром, состоящим из аргументов данной функции. По окончании обработки определения функции кадр этой функции удаляется из окружения. Таким образом, происходит симуляция поведения окружения при выполнении программы на этапе компиляции, чтобы заранее определить положения переменных в окружении и избавиться от поиска переменных по их имени во время выполнения программы.

\paragraph{Константы и переменные}

Константы переносятся в искомое дерево без изменений в форму CONST. Пример:

\begin{figure}[H]
\begin{lstlisting}[language=Lisp]
(CONST 86)
\end{lstlisting}
\end{figure}

При анализе переменной определяется её тип. Должно быть определено 3 типа переменных:

\begin{enumerate}
\item Глобальные переменные. Каждая глобальная переменная заносится в отдельный список *global-vars*. Обращение к таким переменным происходит по индексу переменной в этом списке с помощью формы GLOBAL-REF:

\begin{figure}[H]
\begin{lstlisting}[language=Lisp]
(GLOBAL-REF 0) ; получить значение глобальной переменной со смещением 0
\end{lstlisting}
\end{figure}

В начале компиляции в список глобальных переменных добавляются символы T и NIL. С учётом этого, память для глобальных переменных виртуальный машины должна иметь размер как минимум в 2 элемента, первый из которых должен по умолчанию иметь значение символа T, второй -- значение пустого списка. Таким образом, обращение к символу T принимает вид (GLOBAL-REF 0), а к символу NIL или пустому списку -- (GLOBAL-REF 1).

\item Локальные переменные -- переменные, определённые в текущем кадре активации. Эти переменные заносятся в текущий кадр лексического окружения, и обращение к ним происходит по смещению в этом кадре:

\begin{figure}[H]
\begin{lstlisting}[language=Lisp]
(LOCAL-REF 0) ; получить значение локальной переменной со смещением в кадре 0
\end{lstlisting}
\end{figure}

\item Свободные переменные -- переменные, определённые в одном из более верхних кадров активации по отношению к текущему. Эти переменные также заносятся в лексическое окружение, и обращение к ним происходит с помощью указания смещения относительно текущего кадра и смещения внутри кадра:

\begin{figure}[H]
\begin{lstlisting}[language=Lisp]
(DEEP-REF 1 0) ; получить значение свободной переменной со смещением кадра 1 и со смещением в кадре 0
\end{lstlisting}
\end{figure}

\end{enumerate}

Таким образом, в результате статического анализа все обращения к переменным будут заменены на смещения, благодаря чему не придётся проводить их поиск во время выполнения программы.

\paragraph{progn}

Последовательность выражений progn преобразуется в список, в котором первый элемент -- символ SEQ, а остальные -- операции промежуточной формы, соответствующие подвыражениям исходного progn в том же порядке появления. Для progn с единственным подвыражением вместо формы SEQ генерируется только операция для вычисления подвыражения. Пустой progn преобразуется в вычисление пустого списка. Пример:

\begin{figure}[H]
\begin{lstlisting}[language=Lisp]
(progn 1 2 3) -> (SEQ (CONST 1) (CONST 2) (CONST 3))
(progn 1) -> (CONST 1)
(progn) -> (CONST ())
\end{lstlisting}
\end{figure}

\paragraph{if}

Условная форма if трансформируется в список, первый элемент которого -- символ ALTER, второй -- операция для вычисления условия, третий -- операция для вычисления по первой ветке, четвёртый -- по второй ветке. Пример:

\begin{figure}[H]
\begin{lstlisting}[language=Lisp]
(if t 1 2) -> (ALTER (GLOBAL-REF 0) (CONST 1) (CONST 2))
\end{lstlisting}
\end{figure}

\paragraph{setq}

При анализе присваивания setq определяется тип переменной, затем генерируется соответствующая операция присваивания (GLOBAL-SET для глобальных переменных, LOCAL-SET -- для локальных и DEEP-SET -- для свободных) вместе с рекурсивной генерацией подвыражения. Пример:

\begin{figure}[H]
\begin{lstlisting}[language=Lisp]
(setq x 5) -> (GLOBAL-SET 2 (CONST 5))
\end{lstlisting}
\end{figure}

\paragraph{Функции и замыкания}

Пользовательские функции с фиксированным количеством аргументов обрабатываются аналогично лямбда-выражениям. Имя функции, смещение кадра окружения, количество аргументов сохраняется в таблицу *fix-functions* для последующей проверки при вызове. Для лямбда-функций имена генерируются автоматически.

Функции с переменным числом аргументов сохраняются в таблицу *nary-functions*, в которой хранится имя функции, смещение кадра окружения и число фиксированных аргументов (остальные аргументы добавляются в виде списка).

При обработке формы labels текущее лексическое окружение расширяется определениями локальных функций. Сохраняется имя самой локальной функции и автоматически сгенерированное имя функции. Дальше генерируется код для этих функций, после чего окружение восстанавливается к предыдущему состоянию.

Примитивы с фиксированным числом аргументов хранятся в таблице *fix-primitives*. Хранится имя примитива и число аргументов.

Примитивы с переменным числом аргументов хранятся в таблице *nary-primitives*. Хранится имя примитива и число фиксированных аргументов.

При компиляции применения функции, определяется тип функции: глобальная, лямбда или примитив. Затем происходит проверка числа аргументов. Для лямбда-функций перед вызовом вставляется скомпилированное тело функции. Перед вызовом функции активируется замкнутое окружение функции, создаётся кадр активации с заранее вычисленным размером, вычисляются слева направо аргументы и заносятся в текущий кадр. При вызове примитивов кадр активации не создается.

Для компиляции тела функций используется форма LABEL, которая представляет собой именованную метку, на которую можно совершить переход, что можно использовать для выполнения функций в байт-коде. Для генерации определения функций можно использовать имя функции в качестве названия метки.

Для перехода на метку функции с фиксированным числом аргументов используется форма FIX-CALL, для функции с переменным числом аргументов -- NARY-CALL. Форма FIX-CALL содержит название метки, номер кадра окружения, в котором была определена функция и относительно которой нужно создавать новый кадр, и список операций, каждая из которых вычисляет аргумент функции, в порядке передачи аргументов функции. В форме NARY-CALL дополнительно содержится количество фиксированных аргументов функции.

Для возврата из функции обратно в место после вызова функции используется форма RETURN.

Пример промежуточной формы определения функции через defun и её вызова:

\begin{figure}[H]
\begin{lstlisting}[language=Lisp]
(defun f (x)
  x)
->
(SEQ
  (LABEL F
    (SEQ
      (LOCAL-REF 0)
      (RETURN)))
  (FIX-CALL F 0 ((CONST 5)))) ; переход на метку функции с фиксированным числом аргументов, с переходом на кадр под номером 0 и со списоком аргументов, состоящих из константы 5
\end{lstlisting}
\end{figure}

\paragraph{tagbody}

Форма tagbody преобразуется в форму SEQ с заменой символов на формы LABEL и формы go на GOTO -- безусловный переход на метку. Пример:

\begin{figure}[H]
\begin{lstlisting}[language=Lisp]
(tagbody
  (go skip)
  1
  2
  skip
  3)
->
(SEQ
  (GOTO G1)
  (CONST 1)
  (CONST 2)
  (LABEL G1
    (CONST 3)))
\end{lstlisting}
\end{figure}

\paragraph{catch/throw}

Форма catch преобразуется в аналогичную форму CATCH, и таким же образом форма throw превращается в форму THROW. Пример:

\begin{figure}[H]
\begin{lstlisting}[language=Lisp]
(catch 'test
  (throw 'test 5))
->
(CATCH (CONST TEST) (THROW (CONST TEST) (CONST 5)))
\end{lstlisting}
\end{figure}

\paragraph{Операции после первой фазы}

\renewcommand{\arraystretch}{0.8} % уменьшение расстояний до сетки таблицы
\begin{xltabular}{\textwidth}{|p{6cm}|X|}
\caption{Операции после первой фазы\label{commands:table}}\\
\hline \centrow \setlength{\baselineskip}{0.7\baselineskip} Операция & \centrow \setlength{\baselineskip}{0.7\baselineskip} Описание операции \\
\hline \centrow 1 & \centrow 2\\
\endfirsthead
\caption*{Продолжение таблицы \ref{commands:table}}\\
\hline \centrow 1 & \centrow 2\\ \hline
\finishhead
\hline NOP & Нет операции \\
\hline CONST expr & Выражение константа. \\
\hline LOCAL-REF j & Обращение к j-той локальной переменной. \\
\hline GLOBAL-REF i & Обращение к i-той глобальной переменной. \\
\hline DEEP-REF i j & Обращение к дальней свободной переменной. i -  смещение кадра активации, j - позиция в кадре.\\
\hline SEQ <expr list> & Последовательность вычислений списка выражений. \\
\hline ALTER cond true false & Условный оператор с выражениями условия,  истины и лжи.\\
\hline LOCAL-SET j expr & Присваивание выражения в j-тую локальную  переменную.\\
\hline GLOBAL-SET i expr & Присваивание выражения в i-тую глобальную  переменную.\\
\hline DEEP-SET i j expr & Присваивание выражения в дальнюю свободную  переменную. i - смещение кадра активации, j - позиция в кадре.\\
\hline LABEL lab expr & Выражение с меткой. Выражение может быть пустым. \\
\hline FIX-LET num args expr & let форма, num - число аргументов, args  - список выражений аргументов, expr - тело lambda\\
\hline FIX-CLOSURE name body & Замыкание функции с именем name и телом  body.\\
\hline PRIM-CLOSURE name & Замыкание примитива (фиксированное число аргументов) с именем name.\\
\hline NPRIM-CLOSURE name & Замыкание примитива (переменное число аргументов) с именем name.\\
\hline RETURN & Выход из функции.\\
\hline FIX-CALL name env args & Вызов функции с именем name, номером окружения ofs и списком выражений аргументов args (фиксированное число аргументов).\\
\hline NARY-CALL name num env args & Вызов функции с именем name, номером окружения ofs и списком выражений аргументов args (переменное число аргументов). num - число фиксированных аргументов.\\
\hline FIX-PRIM name args & Вызов примитива (фиксированное число аргументов) с именем name и списком выражений аргументов args.\\
\hline NARY-PRIM name num args & Вызов примитива (переменное число аргументов) с именем name, числом фиксированных аргументов num и списком выражений аргументов args.\\
\hline GOTO name & Переход на метку name.\\
\hline CATCH tag\_exp body\_expr & Блок catch с выражением метки tag\_expr и множеством выражений body\_expr..\\
\hline THROW tag\_expr ret\_expr & Выход из блока catch с выражением метки tag\_expr и результатом вычислений ret\_expr.\\
\end{xltabular}
\renewcommand{\arraystretch}{1.0} % восстановление сетки